\begin{papercard}{tornberg2024}
\cf{Stages} \stagetag{L} (the whole paper), \stagetag{C} (``prompt codebook'' joins annotation guidelines and LLM settings), \stagetag{I} (validation and human--LLM reliability), \stagetag{R} (what to report). A 13-page best-practice essay in \emph{Sociologica} (2024) with nine points for using large language models (LLMs) for text annotation in social research. It is not an empirical study; its recommendations are drawn from the author's projects, prior literature and discussions, and are the most checklist-like source for stage L.

\cf{Position} LLM annotation is ``something of an academic Wild West'': no benchmarks for authors or reviewers, risk of low-quality research \pc{tornberg2024}{2}{LLM-based text annotation has become something of an academic Wild West}. LLMs are subjective, flawed and potentially biased, but so are human coders, and ``in conventional coding procedures, such issues are managed by organizing coding in rigorous processes that identify disagreements, validate the reliability, and make transparent the management of subjectivity''; LLMs should be integrated into such a process, neither used uncritically nor rejected \pc{tornberg2024}{2}{In conventional coding procedures, such issues are managed by organizing coding in rigorous processes that identify disagreements, validate the reliability, and make transparent the management of subjectivity}. Nine points are covered: model choice, systematic coding procedure, prompt codebook, validation, prompt engineering, parameters, ethics and law, stochasticity, training-data contamination (the paper's own numbering of sections differs slightly) \pc{tornberg2024}{3}{We will cover the following nine points}.

\cf{1. Choose an appropriate model} Platform models change without notice and are deprecated, so results become hard to reproduce \pc{tornberg2024}{3}{proprietary models such as ChatGPT have been shown to change over time without notice}; their training data is unknown (low transparency); and they are problematic for confidential data. Six selection factors: reproducibility (``use a fixed version of the LLM throughout the project, document the version, and ensure that the model will be available for future use'') \pc{tornberg2024}{3}{use a fixed version of the LLM throughout the project, document the version, and ensure that the model will be available for future use}; ethics and legality (privacy, no storage of research data); transparency; culture and language (many models are English and US centric); scalability (cost, rate limits; possibly distil into a smaller semi-supervised model); complexity of the task. Best practice: an open-source model with publicly known training data, ideally hosted on one's own infrastructure; assess whether results replicate across several models \pc{tornberg2024}{4}{In general, best practice is to use an open-source model for which the training data is publicly known}; avoid chat web interfaces because they offer no parameter control or version control and may keep the data for training \pc{tornberg2024}{5}{Using LLMs for annotation through their web interface should in general be avoided}. Guardrails can make models refuse controversial content or give low-quality answers \pc{tornberg2024}{4}{the models may provide low-quality responses, or refuse to respond altogether}.

\cf{2. Follow a systematic coding procedure} Annotation involves operationalising concepts; subjectivity is inevitable and should be made explicit. LLMs code one text at a time and do not remember earlier items, unlike humans whose definitions drift. Annotation is iterative: start with texts, discuss discrepancies, refine guidelines, repeat; with an LLM, ``the development of the prompt is simply brought into this loop'' \pc{tornberg2024}{6}{When coding with an LLM, the development of the prompt is simply brought into this loop}. Procedure: (1) define the concept explicitly before consulting the LLM, with separate but similar instructions for humans and for the model; (2) human coders code a small representative set which the LLM also annotates; (3) examine disagreements between coders and between coders and LLM, ask the model to justify cases (but ``coders risk being overly swayed by the model's interpretations, as the models can provide highly convincing explanations'') \pc{tornberg2024}{7}{coders risk being overly swayed by the model’s interpretations, as the models can provide highly convincing explanations}; (4) refine codebook and/or prompt, noting that ``the human coders should not be considered ground truth'' \pc{tornberg2024}{7}{The human coders should not be considered ground truth: you may find that the LLM’s interpretation was superior to the human coder}; (5) repeat; (6) validate on a separate set. Compute reliability among humans and between humans and LLM and aim for the LLM to approach the human level \pc{tornberg2024}{6}{Ideally, the LLM should approach the reliability achieved among the human coders}. If zero-shot is inadequate use few-shot, then consider fine-tuning.

\cf{3. Prompt codebook} A codebook with clear category definitions, examples and rules for ambiguous cases, plus a section giving the exact LLM prompt, model and parameters, so the coding is fully reproducible. Prompts are model-specific: ``applying the same prompt to a different LLM may produce different results, even with models of similar parameter size'' \pc{tornberg2024}{7}{applying the same prompt to a different LLM may produce different results, even with models of similar parameter size}; if the model is fine-tuned, publish the training data.

\cf{4. Engineer prompts} Structured prompt = context (role + background), question, constraints (output format), as in the example of a misinformation annotation with JSON output and an ``Uncertain'' option \pc{tornberg2024}{8}{An annotation prompt should contain the following elements: context, question, and constraints}. Order: context $\to$ instructions $\to$ constraints, because ``recent and repeated text in the prompts has the most effect''. Enumerate categorical options (alphabetical, line-separated); give an explicit ``I don't know'' option to reduce stochastic answers \pc{tornberg2024}{9}{Provide an option for the LLM to respond if it is uncertain about the correct answer}; use lists for complex instructions; use JSON for multiple fields; use an LLM to improve the prompt iteratively; balance brevity and specificity (``performance can fall with longer prompts'', and cost rises); chain-of-thought where reasoning helps; system vs.\ user messages; few-shot examples as user/assistant turns. Prompt engineering is a social-science skill and a form of qualitative conceptual work.

\cf{5. Validate the model} Performance depends on dataset and task and cannot be predicted, so ``it is always necessary to carefully validate the models on a task-by-task basis'' \pc{tornberg2024}{9}{it is always necessary to carefully validate the models on a task-by-task basis}. Requirements: (i) validate \emph{after} the prompt is finalised; using validation data to tune the prompt is not acceptable \pc{tornberg2024}{10}{it is not acceptable to use the validation data to improve the prompts, as this may lead to falsely reporting higher precision}; (ii) sufficient size, at least 20--30 examples per category as a practical minimum, more for imbalanced or high-stakes work, ideally by power analysis \pc{tornberg2024}{10}{The practical minimum is to have at least 20–30 samples of each category for a basic level of confidence in the performance metrics}; (iii) appropriate metrics: accuracy can mislead (90\,\% accuracy by predicting the majority class) \pc{tornberg2024}{10}{if, for instance, one of your categories represents 90\% of the population, then a model that classifies everything as belonging to that category will achieve a seemingly impressive accuracy of 90\%}, so use F1 with precision and recall, weighted F1, ROC-AUC or Cohen's $\kappa$, argue the choice, and combine metrics; (iv) compare with human performance, since some tasks are intrinsically hard (human inter-coder reliability is a benchmark) \pc{tornberg2024}{10}{Calculating the performance of human coders, using e.g., an inter-coder reliability score, can provide a useful benchmark}; (v) validate subsets (languages, cultures) separately; (vi) examine and explain failures, since a 30\,\% error rate on a 10\,\% minority is only a 3\,\% overall failure rate that hides bias \pc{tornberg2024}{10}{Say, for instance, that 10\% of the data describes a particular minority, and that 30\% of these are misclassified due to model bias}.

\cf{6. Parameters} Max length, temperature (low, e.g.\ 0, for annotation), Top-P (0.2--0.4 suggested; temperature must be above 0 when Top-P is used), Top-K. ``Your parameters must always be explicitly specified, even if they are the default parameters'' \pc{tornberg2024}{11}{Your parameters must always be explicitly specified}.

\cf{7. Ethics and law} Informed consent for LLM analysis of personal data, data-processing agreements with third-party providers, changing expectations of privacy for public posts (users did not anticipate LLM inference), anonymisation or pseudonymisation before sending data, data minimisation, storage and transfer outside the EU/EEA, copyright and terms-of-service. ``Never send proprietary, secret, or confidential data to an API or web interface without careful consideration'' \pc{tornberg2024}{12}{Never send proprietary, secret, or confidential data to an API or web interface without careful consideration of the ethical and legal implications}.

\cf{8. Stochasticity and prompt stability} Minor prompt details (capitalisation, punctuation, order) can move accuracy from state-of-the-art to near chance \pc{tornberg2024}{12}{even minor details in the prompt}. Check whether the same prompt gives the same result across repeated runs and whether paraphrased prompts agree; create several paraphrases, run on a subset, and compute Krippendorff's $\alpha$ among the outputs (library by Barrie et al.) \pc{tornberg2024}{13}{create several paraphrases of the prompt and run the analysis for a subset of the data}.

\cf{9. Training-data contamination} Because the LLM may have seen the task data, classical train--test separation is less meaningful; validation concerns task capacity, and it matters only that the validation data resembles what the model will be run on. Problems arise if only part of the data is in training data (validate separately) or if validation labels are public (e.g.\ Manifesto Project): ``avoid using publicly available databases as validation data''; ``the burden of evidence is on the validator'' \pc{tornberg2024}{13}{you should try to avoid using publicly available databases as validation data, as they may be in the model’s training data}.

\cf{Limits (card)} Opinion-based best practice from one social-science author; few quantitative benchmarks; the 20--30 per category figure is a heuristic; the advice to prefer open, self-hosted models may conflict with quality needs; ethics section is general, not specific to health-app screenshots or policies; examples are political text; field moving fast, so details such as parameter ranges are provisional.

\cf{Use for this thesis} (1) A ready checklist for stage L and for the thesis' methods paragraph: model name and version fixed and documented; temperature, top-$p$, max length stated; prompt and codebook archived together; separate tuning and validation sets; metrics argued and reported per code; paraphrase stability measured with $\alpha$; human--human reliability reported as the benchmark. (2) Ethics checklist for the LLM stage: no sending of personal data from interviews or participants to a cloud model without a data-processing agreement; policy texts and public store pages are lower risk but still need a note on terms of use. (3) The explicit warning that the human reference is not ground truth supports reporting human--human and human--LLM agreement symmetrically and resolving disagreements by documented adjudication, not by deferring to the model. (4) Treat as procedural guidance, not as evidence that LLM annotation is valid for privacy codes; pair with empirical work (Gilardi, Pangakis, Baumann) and with Rodriguez et al.\ for the policy domain.
\end{papercard}
